\scanpage{162}
\mspage{1 - 38 -}
\margin{n° 413}

\uline{Exemple} 1.- Soit $\vartheta$ l'application $x\longrightarrow x^{-1}$ de $G$ dans lui-même : c'est un isomorphisme de $G$ sur la variété de groupe opposée $G^o$. Si $t\in T^{(r)}(G)$, on désigne encore par $\check t$ l'image de $t$ par $\vartheta$. Avec cette notation, on a :
\[
(23)\qquad (t*t')^\vee=\check t'*\check t
\]
quelles que soient les distributions ponctuelles $t$ et $t'$ d'ordre $\leq s$ et $s'$ sur $G$, avec $s+s'\leq r$.

\uline{Exemple} 2.- Soit $H$ une sous-variété de groupe de $G$. L'injection canonique de $H$ dans $G$ permet d'identifier $\mathcal T^{(r)}(H)$ à un sous-espace vectoriel de $\mathcal T^{(r)}(G)$ et la prop. 20, (ii), montre que cette identification est compatible avec le produit de convolution.

Soient maintenant $H_i$ $(i=1,2)$ deux sous-variétés de groupe de $G$ tels que tout élément de $H_1$ commute avec tout élément de $H_2$ ; soit $t_i\in T^{(s_i)}(H_i)$, avec $s_1+s_2\leq r$. On a alors $t_1*t_2=t_2*t_1$ car l'une comme l'autre de ces deux distributions est l'image de $t_1\otimes t_2$ par l'application $(h_1,h_2)\longrightarrow h_1h_2=h_2h_1$ de $H_1\times H_2$ dans $G$.

Si en particulier $G$ est \uline{commutatif}, alors $\mathcal T^{(r)}(G)$ est \uline{commutative}.

Si $r\geq\infty$, la prop. 20 montre que le produit de convolution fait de l'espace vectoriel $\mathcal T^{(r)}(G)$ une \uline{algèbre associative} et que $\varphi_*$ est un \uline{homomorphisme d'algèbres} pour tout morphisme $\varphi$ de classe $C^\infty$ de $G$ dans $H$. En particulier, $\mathcal T^{(r)}(G^o)$ est l'algèbre \uline{opposée} à $\mathcal T^{(r)}(G)$.

Le sous-espace $T_{\unsure{o|0}}^{(r)}(G)$ de $\mathcal T^{(r)}(G)$ est \uline{stable} pour la loi de composition non partout définie $*$. \uline{Nous le noterons encore} $U(G)$. Pour $r\geq\infty$, c'est une algèbre associative que nous étudierons plus en détail au § 3. Nous noterons
\margin{161}
